IVF围绕粗质心分区，列表保存归属于各质心的向量。
查询选$n_{\rm probe}$个列表，得到含$n_c$个候选的路由全集$U(q)$。
少访问列表可以省工，却可能在精细评分前排除真近邻。
质心用于路由，其距离选择与选择数据库向量是两个不同问题。

PQ把向量分成$M$个坐标子空间，每个子空间用小型学习码本的标签表示\cite{pq}。
若$\widehat x_m=\mathcal C_m[z_m]$，ADC保持查询不量化，计算
\begin{equation}
 \|q-\widehat x\|^2=\sum_{m=0}^{M-1}\|q_m-\mathcal C_m[z_m]\|^2.
 \label{eq:2-adc}
\end{equation}
查询表预计算子空间距离，扫描执行code load、查表与累加。
乘积存在于子空间之间，学习子码本\emph{内部}的质心未必还有坐标乘积。
这决定第5章的精确内部拆表是否成立。

这里的“残差”有两种含义。IVF-PQ常编码$x-\mu_{\ell(x)}$，粗质心参与重建；
JHQ用IVF路由，却相对旋转后的primary量化重建定义残差。
两者依赖不同；给JHQ残差重建再加IVF质心，会加入编码器未减去的量。
原PQ/ADC框架及GPU实现也支持扩大候选后进行更精确的refinement
\cite{johnson_gpu,faiss_refine}。

在至少有$k$个有效路由候选的分析范围内，JHQ保留有效primary集合$C_\alpha(q)$，容量上限
$c(\alpha)=\max\{k,\lceil\alpha k\rceil\}$，精排返回$\mathcal R_\alpha(q)$：
\begin{equation}
 \mathcal R_\alpha(q)\subseteq C_\alpha(q)\subseteq U(q)\subseteq\mathcal X.
 \label{eq:2-cascade}
\end{equation}
与GT求交得到各阶段覆盖对最终recall的上界，却不能定量归因各阶段损失；
后者需要固定候选全集的反事实评分。精排可以改善保留向量之间的次序，无法恢复被排除ID。

\par\medskip\noindent\begin{minipage}{\linewidth}
\centering\captionsetup{hypcap=false}
\resizebox{\linewidth}{!}{\begin{tikzpicture}[
 font=\small\linespread{1}\selectfont,
 box/.style={draw=ThesisNavy!65,rounded corners=2pt,align=center,
 text width=2.65cm,minimum height=1.8cm,inner sep=4pt},
 arr/.style={-{Latex[length=2mm]},draw=ThesisNavy!80,thick}]
\node[box] (x) at (0,0) {Database $\mathcal X$\\$N$ vectors};
\node[box] (u) at (3.6,0) {IVF lists\\$U(q)$: $n_c$ vectors};
\node[box,fill=ThesisNavy!7] (c) at (7.2,0) {Primary top-$c$\\$C_\alpha(q)$\\up to $c$ vectors};
\node[box,fill=ThesisNavy!7] (r) at (10.8,0) {Refined top-$k$\\$\mathcal R_\alpha(q)$};
\draw[arr] (x) -- node[above,font=\footnotesize] {IVF} (u);
\draw[arr] (u) -- node[above,font=\footnotesize] {$D_{\rm P}$} (c);
\draw[arr] (c) -- node[above,font=\footnotesize] {$D_{\rm H}$} (r);
\node[align=center,text width=6.4cm] at (2.6,-2.0)
 {Query $q'=\Pi q$ stays FP32.\\Primary codes score the broad scan.};
\node[align=center,text width=5.8cm] at (9.7,-2.0)
 {Residual $r=\Pi x-\widehat p$.\\Survivors decode $\widehat p+\widehat r$.};
\end{tikzpicture}
}
\captionof{figure}[候选与精度的层级流水线]{候选与精度级联。IVF选列表，primary选向量，只有survivor读残差。
层次结构是继承的；图定义逻辑集合与评分精度，不表示实测recall损失或GPU内核边界。}
\label{fig:2-cascade}
\end{minipage}\par\medskip

忽略路由与选择，广泛查表评分需要$O(n_cM)$次lookup，残差精排需要$O(cd)$坐标工作。
这一不对称推动选择性精度，却不能预测延迟，还必须考虑内存、算术、并行及候选选择。
增加probe改变$U$，增加$\alpha$改变该全集内的survivor；把二者都称搜索深度会掩盖差异。
